\section{Teoría elemental de números}

\begin{ejercicio}
    Para ningún número natural $n$ es $n^2 + 1$ divisible por 7.
\end{ejercicio}

\begin{proof}
    $n^2 + 1$ es divisible entre 7 si y solo si $n^2 = 7k + 6$ para algún entero $k$, es decir,
    si y solo si $(n^2) \MOD 7 = 6$.
    Basta, por tanto, con ver que $n^2 \MOD 7 \neq 6$ para todo $n$.
    Observamos que:
    \[ n \MOD 7 = (n \MOD 7) \MOD 7 \]
    y por tanto, utilizando la propiedad 2 del módulo,
    \[ n^2 \MOD 7 = {(n \MOD 7)}^2 \MOD 7 \]
    ya que $0 \le n \MOD 7 < 7$, será suficiente comprobar los 7 valores posibles de $n \MOD 7$.
    \begin{align*}
        0^2 \MOD 7 &= 0, & 1^2 \MOD 7 &= 1, & 2^2 \MOD 7 &= 4, & 3^2 \MOD 7 &= 2, \\
        4^2 \MOD 7 &= 2, & 5^2 \MOD 7 &= 4, & 6^2 \MOD 7 &= 1.
    \end{align*}
    Ningún resultado es 6, luego $n^2 \MOD 7 \neq 6$ para todo $n$.
\end{proof}

\begin{definicion}\label{def:mcd}
    Sean $a$ y $b$ enteros. Un entero $d \neq 0$ es un divisor común de $a$ y $b$ si
    $d \mid a$ y $d \mid b$. Un divisor común de $a$ y $b$ se llama
    \emph{máximo común divisor} de $a$ y $b$ si $d > 0$ y cada común divisor de
    $a$ y $b$ divide también a $d$. Al máximo común divisor de $a$ y $b$
    le designaremos por $\mcd(a, b)$.
    
    En el caso que $a = b = 0$, entonces $\mcd(0, 0) = 0$.
\end{definicion}

\begin{teorema}[Identidad de Bézout]\label{teo:bezout}
    Sean $a$ y $b$ enteros distintos de $0$. Entonces existe un único $d$ máximo
    común divisor de $a$ y $b$. Además $d$ es el entero positivo más pequeño que
    puede expresarse en la forma $ax + by$ donde $x$ e $y$ son números enteros.
\end{teorema}

\begin{proof}
    Probemos que existe $d > 0$ que divide a $a$ y $b$, y
    a su vez es divisible por cada común divisor de $a$ y $b$.

    Tomando cualquier divisor común $c$,
    por la propiedad 6 del apartado 1.3 del texto se cumple que $c \mid (ax + by)$
    para cualquier par de enteros $x$ e $y$.

    Ahora probemos la existencia del menor elemento que tenga la forma $ax + by$.
    Sea $S$ el conjunto de números $z$ que cumplen la propiedad $z = ax + by$ con $z > 0$.
    $S$ es un conjunto no vacío dado que, para $a > 0$ tomando el par $x = 1$ y $y = 0$
    resultaría en $z = a$, y para $a < 0$ tomando el par $x = -1$ y $y = 0$ resultaría en $z = -a$.
    Debido a que no es el conjunto vacío y está acotado inferiormente, por el principio
    de la buena ordenación existe un elemento mínimo de $S$, lo denotaremos por $d = ax_0 + by_0$.

    Utilizando el algoritmo de la división: $a = dq + r$, donde $0 \le r < d$.
    Procedamos por reducción al absurdo, supongamos $r > 0$.
    Podemos reescribir $r$ de la siguiente manera:
    \begin{align*}
        r &= a - dq = a - (ax_0 + by_0)q \\
          &= a(1 - x_0q) + b(-y_0q)
    \end{align*}
    y como $(1 - x_0q)$ y $(-y_0q)$ son enteros, $r \in S$. Pero $r < d$ llegando a una contradicción
    ya que $d$ es el menor elemento de $S$, por lo tanto $r = 0$, siendo así $d$ un divisor de $a$.

    Se demuestra la afirmación $d \mid b$ de forma análoga.
    Como $d$ es de la forma $ax + by$ todo divisor común a $a$ y $b$ divide a $d$.
    Por tanto, $d$ es un máximo común divisor de $a$ y $b$, lo que demuestra su existencia.

    Sea $d'$ otro máximo común divisor de $a$ y $b$.
    Como ambos son máximo común divisores, $d, d' > 0$.
    Como $d'$ es divisor común y $d$ es máximo común divisor, $d' \mid d$.
    Y recíprocamente $d \mid d'$.
    Cumplen las condiciones necesarias para aplicar la propiedad 7
    del apartado 1.3 del texto: $d \le d'$ y $d' \le d$.
    Por lo tanto $d = d'$.

    Así, el máximo común divisor es único y coincide con $d$,
    que por construcción es el menor entero positivo de la forma $ax + by$.
\end{proof}


\begin{lema}[Lema de Euclides]\label{lem:euclides}
    Sean $a$, $b$ y $c$ números enteros. Supongamos que $a$ y $c$ son primos entre
    sí y que $c \mid ab$. Entonces $c \mid b$.
\end{lema}

\begin{proof}
    Si $a = 0$, entonces $1 = \mcd(a, c) = \mcd(0, c) = |c|$,
    luego $c = \pm1$ y $c \mid b$.

    Supongamos ahora $a \ne 0$. Como $c \mid ab$, se tiene $c \ne 0$.
    Por tanto, podemos aplicar el Teorema~\ref{teo:bezout}: existen $s$ y $t$ enteros
    tales que $as + ct = 1$. Multiplicando ambos miembros de la igualdad por $b$ obtenemos:
    \[ abs + cbt = b. \]

    Como $c \mid ab$, por definición de divisibilidad existe un entero $k$ tal que
    $ab = ck$. Sustituyendo $ab$ en la igualdad anterior:
    \begin{align*}
        cks + cbt &= b \\
        c(ks + bt) &= b.
    \end{align*}
    Dado que $ks + bt$ es un número entero, $c$ divide a $b$.
\end{proof}